\chapter{Conclusion and Future Work}
\label{ch:conclusion}

\section{Conclusion}

This project set out to test whether an SSH honeypot that chooses its responses can capture more
attacker behaviour than one that does not, and to measure the benefit rigorously.

We built \textbf{CruxTrap}, a complete five-service system. It consists of a safe
medium-interaction honeypot with a per-session emulated Ubuntu server and over 200 emulated
command names, an online reinforcement-learning response policy, an at-least-once event pipeline
through Redis Streams into PostgreSQL, and an analyst dashboard with session replay and a view of
the learned policy. The system is containerised with defence-in-depth hardening. Its key safety
and availability properties were verified on the running deployment: an unprivileged, read-only,
capability-free honeypot that cannot reach the database, and continued attack capture through
Redis, pipeline and database outages.

We formalised the honeypot's problem as an MDP whose reward measures captured intelligence, and
identified that it is partially observable because the attacker's type is hidden. For evaluation
we built a generative attacker model and an \emph{exact} evaluator. The evaluator computes each
policy's expected return, engagement and intelligence by solving the Bellman equations on a
330-state augmented chain; it also gives an oracle upper bound and the best policy available to
the honeypot's own state.

The main findings are:
\begin{enumerate}
  \item Q-learning beats a static honeypot with high confidence ($5.686$ against $5.542$,
    $p = 5.8 \times 10^{-6}$) and recovers 60\% of the achievable improvement. It does this by
    discovering, without supervision, a single robust deception: serving decoy credentials to
    attackers who are searching for credentials.
  \item This holds only with a step size that satisfies the Robbins--Monro conditions. With the
    conventional constant learning rate the agent did worse than static, because the estimates
    never settle and noise decides the greedy action. An ablation confirmed the explanation.
  \item The benefit is modest and depends on the population: it is negligible against scripted
    bots and about 5\% against human-heavy traffic. Adaptation pays where attackers adapt.
  \item Exploration while live costs about as much as adaptation gains, so a deployment should
    reduce exploration once the policy has settled.
\end{enumerate}

The broader lesson is methodological. Deploying an adaptive honeypot and reporting that
attackers stayed longer does not establish that adaptation helped. A model with exact evaluation
makes it possible to ask how much better a policy is, compared with what is achievable, and
against whom.

\section{Future Work}

\begin{itemize}
  \item \textbf{Real-traffic study.} Deploy the honeypot on a Linux host with a public address.
    Fit the attacker model to sessions captured under the static policy, then compare the learned
    and static policies on live traffic with an interleaved (A/B) design.
  \item \textbf{Richer state.} Add the previous action, a short command history, or a belief over
    the attacker type estimated from inter-command timing. The posterior in
    Figure~\ref{fig:posterior} suggests that even a coarse type estimate is valuable. Beyond a few
    hundred states, linear function approximation or a small DQN becomes appropriate.
  \item \textbf{Exploration and bias.} Use GLIE schedules for deployment, optimistic
    initialisation or double Q-learning for the late-state bias, and off-policy evaluation from
    logged data.
  \item \textbf{Durability.} Use a disk-backed spool in the emitter so that events survive a
    simultaneous Redis outage and honeypot crash.
  \item \textbf{Capture.} Add SFTP/SCP file capture, sandboxed analysis of downloaded payloads, and
    threat-intelligence enrichment of indicators.
  \item \textbf{Operations.} Add dashboard authentication and HTTPS, alerting, SIEM export, and
    orchestration (for example Kubernetes) for multi-sensor deployments.
\end{itemize}
